\begin{definition}[Algebra conditions for a basis of normal tensors]%
    \label{def:mpdo_bnt_algebra_clause}
    \lean{MPOTensor.BNTAlgebraClause}
    \leanok
    \uses{def:mpdo_bnt_label_coefficient_family, def:mpdo_bnt_label_operator_family, def:mpdo_bnt_label_trace_scalar_family, def:mpdo_diagonal_chi_family, def:mpdo_positive_bnt_label_chi_trace_power_form, def:mpdo_bnt_label_same_length_product_form, def:mpdo_bnt_label_idempotent_coefficient_form}
    Fix structure coefficients $c^{(L)}_{\alpha,\beta,\gamma}$,
    operators $O_L(M_\alpha)$, and scalars $m_\alpha=\tr(\mu_\alpha)$.
    The algebra conditions require a length-independent family of
    positive diagonal matrices $\chi_{\alpha,\beta,\gamma}$ such that, for
    every $L>0$,
    $c^{(L)}_{\alpha,\beta,\gamma}
        =\tr(\chi_{\alpha,\beta,\gamma}^{L})$ and
    \begin{align}
        O_L(M_\alpha)O_L(M_\beta)
         & =\sum_\gamma c^{(L)}_{\alpha,\beta,\gamma}O_L(M_\gamma),
        \notag
    \end{align}
    together with
    \begin{align}
        m_\gamma
         & =\sum_{\alpha,\beta}
        c^{(1)}_{\alpha,\beta,\gamma}m_\alpha m_\beta.
        \notag
    \end{align}
    Thus the vector $m=(m_\alpha)_\alpha$ is an idempotent for the
    multiplication induced by $c^{(1)}$. These are the algebra conditions in
    \cite[Theorem~4.14(ii), labelled IV.13 in the source,
        lines~972--985]{Cirac2016MPDO_arXiv}, without an additional comparison
    with a chosen blocked basis.
\end{definition}

\begin{definition}[Algebra conditions for a vertical canonical decomposition]%
    \label{def:mpdo_bnt_algebra_tensor_clause}
    \lean{MPOTensor.verticalBNTMPO}
    \lean{MPOTensor.verticalBNTOperatorFamily}
    \lean{MPOTensor.verticalBNTTraceScalarFamily}
    \lean{MPOTensor.BNTAlgebraTensorClause}
    \lean{MPOTensor.HasBNTAlgebraTensorClause}
    \leanok
    \uses{def:vertical_cf_mpdo, def:mpdo_bnt_algebra_clause}
    Choose a vertical canonical decomposition
    \begin{align}
        U\widetilde M U^\dagger
         & =\bigoplus_\alpha \mu_\alpha\otimes M_\alpha,
        \notag                                                          \\
        \mu_\alpha
         & =\diag(\omega_{\alpha,0},\ldots,\omega_{\alpha,r_\alpha-1}).
        \notag
    \end{align}
    Here the $\mu_\alpha$ are positive diagonal matrices with nonempty
    multiplicity spaces, and $U$ is a coisometry onto the retained sectors,
    so $UU^\dagger=\Id$. The decomposition also requires exact
    reconstruction:
    \begin{align}
        \widetilde M
         & =U^\dagger
        \left(\bigoplus_\alpha\mu_\alpha\otimes M_\alpha\right)U.
        \notag
    \end{align}
    These tensor identities hold for every vertical physical index.
    In particular, any complementary sector is zero.
    Separating the two factors of the doubled physical index of
    $M_\alpha$ gives an MPO tensor $\widehat M_\alpha$. Define
    \begin{align}
        O_L(M_\alpha)
         & =\operatorname{mpo}(\widehat M_\alpha,L),
        \notag                                       \\
        m_\alpha
         & =\sum_q\omega_{\alpha,q}=\tr(\mu_\alpha).
        \notag
    \end{align}
    The chosen decomposition satisfies the algebra conditions when the
    tensors $M_\alpha$ form a BNT of the vertical tensor and the conditions
    of Definition~\ref{def:mpdo_bnt_algebra_clause} hold for these operators
    and scalars. These are the objects and equations in
    \cite[Theorem~4.14(ii), lines~972--993, and Appendix~C.4,
        lines~2046--2064]{Cirac2016MPDO_arXiv}; a comparison with a two-site
    vertical canonical decomposition is an additional requirement.
\end{definition}

\begin{theorem}[Positivity of the physical trace in a retained vertical sector]%
    \label{thm:mpdo_algebra_clause_terminal_transfer_positivity}
    \lean{MPOTensor.BNTAlgebraTensorClause.%
        physTraceTransfer_verticalBNTMPO_posSemidef}
    \leanok
    \uses{def:mpdo, def:mpdo_bnt_algebra_tensor_clause, def:mpo_physical_trace_transfer, lem:mpdo_sitewise_physical_coordinate_congruence}
    Let $M$ generate a family of MPDOs with physical dimension $d$
    and bond dimension $D$, and choose a vertical canonical decomposition
    satisfying the algebra conditions of
    Definition~\ref{def:mpdo_bnt_algebra_tensor_clause}.
    Let $M_\gamma$ be any retained normal tensor. Separating its
    doubled physical index gives two physical legs of dimension $D$,
    corresponding to the original auxiliary indices of $M$.
    Closing these legs gives the bond matrix
    \begin{align}
        T_\gamma
         & =\sum_{a=0}^{D-1}M_\gamma^{(a,a)}.
        \label{eq:rfp_algebra_terminal_transfer}
    \end{align}
    Then $T_\gamma\succeq0$. This follows from the chosen decomposition
    and positivity of the one-site operator $\rho^{(1)}(M)$; it does
    not require fusion maps or an RFP hypothesis.
\end{theorem}

\begin{proof}\leanok
    \uses{lem:mpdo_sitewise_physical_coordinate_congruence}
    Apply the retained physical coisometry to the positive one-site operator.
    For a positive diagonal coordinate $q$ of the multiplicity matrix
    $\mu_\gamma$, the corresponding principal submatrix is
    $(\mu_\gamma)_{q,q}T_\gamma$. Principal submatrices remain positive
    semidefinite, and $(\mu_\gamma)_{q,q}>0$; division by this scalar proves
    $T_\gamma\succeq0$.
\end{proof}

\begin{lemma}[Physical traces of a nonunitary similarity pair]%
    \label{lem:mpdo_nonunitary_pair_terminal_transfers}
    \lean{MPSTensor.PositiveMinimalRealizationCounterexample.%
        physTraceTransfer_tensor}
    \lean{MPSTensor.PositiveMinimalRealizationCounterexample.%
        physTraceTransfer_gaugedTensor}
    \leanok
    \uses{def:mpo_physical_trace_transfer}
    Let $A$ have letters $\Id,E_{00},E_{01},E_{10}$, let
    $X=\diag(2,1)$, and set $B^v=XA^vX^{-1}$. After separating the doubled
    physical index, closing the physical legs gives
    \begin{align}
        T_A & =\begin{pmatrix}1&0\\1&1\end{pmatrix},
        \notag                                              \\
        T_B & =\begin{pmatrix}1&0\\\tfrac12&1\end{pmatrix}.
        \notag
    \end{align}
\end{lemma}

\begin{proof}\leanok
    Sum the two diagonal physical letters $A^{(0,0)}+A^{(1,1)}$ and then use
    $B^v=XA^vX^{-1}$ with $X=\diag(2,1)$.
\end{proof}

\begin{lemma}[The physical traces of the nonunitary similarity pair are not positive semidefinite]%
    \label{lem:mpdo_nonunitary_pair_terminal_not_positive}
    \lean{MPSTensor.PositiveMinimalRealizationCounterexample.%
        physTraceTransfer_tensor_not_posSemidef}
    \lean{MPSTensor.PositiveMinimalRealizationCounterexample.%
        physTraceTransfer_gaugedTensor_not_posSemidef}
    \leanok
    \uses{lem:mpdo_nonunitary_pair_terminal_transfers}
    Neither $T_A$ nor $T_B$ in
    Lemma~\ref{lem:mpdo_nonunitary_pair_terminal_transfers} is positive
    semidefinite.
\end{lemma}

\begin{proof}\leanok
    \uses{lem:mpdo_nonunitary_pair_terminal_transfers}
    A positive semidefinite matrix is Hermitian. In both matrices the
    $(0,1)$ entry is zero, whereas the $(1,0)$ entry is respectively $1$ and
    $1/2$. Hence neither matrix is Hermitian.
\end{proof}

\begin{theorem}[The nonunitary similarity pair is excluded from retained positive sectors]%
    \label{thm:mpdo_nonunitary_pair_not_retained_positive_sector}
    \lean{MPSTensor.PositiveMinimalRealizationCounterexample.%
        tensor_not_retained_positive_sector}
    \lean{MPSTensor.PositiveMinimalRealizationCounterexample.%
        gaugedTensor_not_retained_positive_sector}
    \leanok
    \uses{def:mpdo_bnt_algebra_tensor_clause, thm:mpdo_algebra_clause_terminal_transfer_positivity, lem:mpdo_nonunitary_pair_terminal_not_positive}
    Neither $A$ nor $B$ can occur, up to the bond-dimension identification,
    as a retained sector of bond dimension two in a vertical canonical
    decomposition of an MPDO tensor satisfying
    Definition~\ref{def:mpdo_bnt_algebra_tensor_clause}.
\end{theorem}

\begin{proof}\leanok
    \uses{thm:mpdo_algebra_clause_terminal_transfer_positivity, lem:mpdo_nonunitary_pair_terminal_not_positive}
    For any such retained sector, closing the physical legs gives a positive
    semidefinite matrix by
    Theorem~\ref{thm:mpdo_algebra_clause_terminal_transfer_positivity}.
    Lemma~\ref{lem:mpdo_nonunitary_pair_terminal_not_positive} gives the
    contradiction for both tensors.
\end{proof}

\begin{lemma}[The nonunitary similarity pair has no positive isometric realization]%
    \label{lem:mpdo_nonunitary_pair_no_positive_isometric_realization}
    \lean{MPSTensor.PositiveMinimalRealizationCounterexample.%
        no_positive_isometric_realization}
    \leanok
    Let $A$ have letters $\Id,E_{00},E_{01},E_{10}$, let
    $X=\diag(2,1)$, and set $B^i=XA^iX^{-1}$. There are no $c>0$ and
    $V\in\MN{2}$ satisfying $V^\dagger V=\Id$ and
    $V^\dagger B^iV=cA^i$ for every letter~$i$.
\end{lemma}

\begin{proof}\leanok
    \uses{thm:normal_tensor_gram_rigidity}
    The identity letter forces $c=1$. The matrix $S=V^\dagger X$ then
    commutes with all four letters of $A$, which span $\MN{2}$, and hence
    $S$ is scalar. Together with $V^\dagger V=\Id$, this would make
    $X^\dagger X=\diag(4,1)$ scalar, a contradiction.
\end{proof}

\begin{definition}[Linear extension on each virtual block]%
    \label{def:mpdo_virtual_perBlockLinearExtension}
    \leanok
    \uses{def:injective, def:same_mpv}
    For injective families $A_k$ and $B_k$ with the same closed matrix
    product vectors, the linear extension on block $k$ is the unique linear map
    $T_k:\MN{D_k}\to\MN{D_k}$ satisfying $T_k(A_k^i)=B_k^i$ for every
    letter~$i$.
\end{definition}

\begin{theorem}[Multiplicativity of the linear extension on each virtual block]%
    \label{thm:mpdo_virtual_perBlockLinearExtension_mul}
    \leanok
    \uses{def:mpdo_virtual_perBlockLinearExtension}
    The linear extension on block $k$ satisfies $T_k(MN)=T_k(M)T_k(N)$ for all
    $M,N\in\MN{D_k}$.
\end{theorem}

\begin{proof}\leanok
    \uses{def:mpdo_virtual_perBlockLinearExtension, thm:linear_ext_mul}
    The letters of $A_k$ span $\MN{D_k}$, and equality of the closed matrix
    product vectors gives the identity on products of spanning letters.
\end{proof}

\begin{theorem}[Bijectivity of the linear extension on each virtual block]%
    \label{thm:mpdo_virtual_perBlockLinearExtension_bijective}
    \leanok
    \uses{def:mpdo_virtual_perBlockLinearExtension}
    The linear extension $T_k$ on block $k$ is bijective.
\end{theorem}

\begin{proof}\leanok
    \uses{thm:mpdo_virtual_perBlockLinearExtension_mul, thm:simplicity, lem:trace_ne_zero}
    Multiplicativity makes the kernel an ideal of the simple matrix algebra.
    The map is nonzero: otherwise every letter of $B_k$ would vanish,
    contradicting the nonvanishing trace pairing for the injective family.
    Thus the kernel is zero, and finite dimensionality gives surjectivity.
\end{proof}

\begin{theorem}[Positive linear extensions on virtual blocks are unitarily implemented]%
    \label{thm:mpdo_positive_perBlockLinearExtension_unitary}
    \leanok
    \uses{def:mpdo_virtual_perBlockLinearExtension}
    If $T_k$ is positive, then there is a unitary matrix $U_k$ such that
    $T_k(M)=U_kMU_k^\dagger$ for every $M\in\MN{D_k}$.
\end{theorem}

\begin{proof}\leanok
    \uses{thm:mpdo_virtual_perBlockLinearExtension_mul, thm:mpdo_virtual_perBlockLinearExtension_bijective}
    The map on each block is a bijective algebra endomorphism. Positivity makes it
    preserve adjoints, so unitary Skolem--Noether implements it by conjugation
    with a unitary matrix.
\end{proof}

\begin{theorem}[Closed matrix product vectors do not force a positive virtual extension]%
    \label{thm:mpdo_closed_mpv_not_force_positive_virtual_extension}
    \lean{MPSTensor.PositiveMinimalRealizationCounterexample.%
        perBlockLinearExtension_not_positive}
    \lean{MPSTensor.PositiveMinimalRealizationCounterexample.tensorFamily}
    \lean{MPSTensor.PositiveMinimalRealizationCounterexample.gaugedTensorFamily}
    \lean{MPSTensor.PositiveMinimalRealizationCounterexample.tensorFamily_isInjective}
    \lean{MPSTensor.PositiveMinimalRealizationCounterexample.gaugedTensor_isInjective}
    \lean{MPSTensor.PositiveMinimalRealizationCounterexample.%
        tensorFamily_sameMPV_gaugedTensorFamily}
    \leanok
    \discussion{5284}
    \uses{def:injective, def:same_mpv, def:mpdo_virtual_perBlockLinearExtension, lem:mpdo_nonunitary_pair_no_positive_isometric_realization}
    Regard the tensors $A$ and $B$ of
    Lemma~\ref{lem:mpdo_nonunitary_pair_no_positive_isometric_realization}
    as one-block families. They are injective and have equal closed matrix
    product vectors at every length, but the linear extension on their
    virtual block is not positive. This concerns the virtual map determined
    by the closed vectors. It does not rule out a comparison of positive
    physical sector realizations, nor does it decide whether an additional
    comparison of chosen vertical decompositions, with specified insertions,
    admits a realization with positive coefficients and identity factors.
\end{theorem}

\begin{proof}\leanok
    \uses{def:mpdo_virtual_perBlockLinearExtension, thm:mpdo_positive_perBlockLinearExtension_unitary, lem:mpdo_nonunitary_pair_no_positive_isometric_realization}
    If the linear extension were positive, its unitary-implementation
    theorem would give a unitary $U$ satisfying
    $B^i=UA^iU^\dagger$ for every letter~$i$.
    Hence $U^\dagger B^iU=A^i$, giving the forbidden realization with
    $c=1$ in
    Lemma~\ref{lem:mpdo_nonunitary_pair_no_positive_isometric_realization}.
\end{proof}
